\documentclass[11pt]{article}
\usepackage{preambulo}
\addbibresource{bibliografia.bib}

\begin{document}

\portada{Evaluación del modelo}

\tableofcontents
\newpage

%============================================================================
\section{Objetivo y protocolo de evaluación}
%============================================================================

El presente documento evalúa los tres clasificadores entrenados sobre el conjunto SUPPORT2
\parencite{support2} para predecir la variable \texttt{hospdead}, documentados en el informe
anterior: regresión logística, árbol de decisión y vecinos más cercanos, acompañados de un modelo
de referencia que predice siempre la clase mayoritaria.

La evaluación se realizó sobre el conjunto de prueba, conformado por 1{,}821 registros que no
intervinieron en el entrenamiento de los modelos ni en la elección de sus hiperparámetros, fijados
por validación cruzada sobre el conjunto de entrenamiento. Debe señalarse que las comparaciones
sobre reducción de dimensionalidad del informe de preprocesamiento sí se evaluaron sobre este
mismo conjunto, de modo que la decisión de conservar las 40 variables no es independiente de él. La estabilidad
de los resultados se verificó adicionalmente mediante validación cruzada sobre el conjunto de
entrenamiento.

El propósito no se limita a determinar cuál modelo obtiene el mejor valor en cada métrica. Un
único número de desempeño admite lecturas erróneas en conjuntos desbalanceados, y una diferencia
entre dos modelos puede deberse al azar de la partición. La evaluación aborda por tanto tres
preguntas: qué miden las métricas empleadas, qué tan estables son los resultados, y si las
diferencias observadas entre modelos son estadísticamente significativas.

%============================================================================
\section{Métricas empleadas}
%============================================================================

\subsection{Por qué la exactitud resulta insuficiente}

La exactitud expresa la proporción de casos clasificados correctamente. Su interpretación directa
la convierte en la métrica de uso más extendido, pero en conjuntos desbalanceados resulta
engañosa.

En el conjunto analizado los fallecimientos representan el 25.9\% de los registros. Un
clasificador que asignara invariablemente la categoría mayoritaria acertaría en el 74.1\% de los
casos sin utilizar ninguna variable de entrada. Ese valor constituye el piso de la métrica, no su
origen: una exactitud de 0.80 no representa un acierto de ocho de cada diez casos sobre un punto
de partida nulo, sino una mejora de seis puntos sobre un modelo que no aprende.

Por ese motivo el modelo de referencia se incorpora a todas las tablas de este documento.

\subsection{Métricas aplicadas}

La evaluación emplea cinco métricas. Las tres últimas se calculan sobre la clase minoritaria, que
corresponde a los pacientes que fallecen y constituye el objeto de interés.

\begin{itemize}
\item \textbf{AUC}: área bajo la curva ROC. Mide la capacidad de ordenar a los pacientes por
      riesgo, con independencia del umbral de decisión. Un valor de 0.5 equivale al azar.
\item \textbf{Exactitud}: proporción de clasificaciones correctas sobre el total.
\item \textbf{Precisión}: de los pacientes señalados como fallecimiento, proporción que
      efectivamente falleció.
\item \textbf{Recall}: de los pacientes que fallecieron, proporción que el modelo señaló.
\item \textbf{F1}: media armónica entre precisión y recall.
\end{itemize}

La distinción entre precisión y recall resulta central en este problema. La primera cuantifica el
costo de las falsas alarmas; la segunda, el de los casos no detectados. Los dos tipos de error no
son intercambiables y su costo relativo depende del contexto de aplicación, que este trabajo no
determina, por lo que resumirlos en un único valor oculta información relevante para la
decisión.

%============================================================================
\section{Resultados sobre el conjunto de prueba}
%============================================================================

La Tabla \ref{tab:metricas} presenta el desempeño de los cuatro modelos.

\begin{table}[H]
\centering
\caption{Métricas sobre el conjunto de prueba. Precisión, recall y F1 corresponden a la clase \textit{fallece}.}
\label{tab:metricas}
\small
\begin{tabular}{p{3.6cm} r r r r r}
\toprule
\textbf{Modelo} & \textbf{AUC} & \textbf{Exactitud} & \textbf{Precisión} & \textbf{Recall} & \textbf{F1} \\
\midrule
Regresión logística & 0.879 & 0.827 & 0.736 & 0.519 & 0.609 \\
Árbol de decisión & 0.847 & 0.807 & 0.701 & 0.447 & 0.546 \\
Vecinos más cercanos & 0.835 & 0.797 & 0.721 & 0.356 & 0.477 \\
Referencia trivial & 0.500 & 0.741 & 0.000 & 0.000 & 0.000 \\
\bottomrule
\end{tabular}
\end{table}


La regresión logística obtiene el mayor valor en las cinco métricas. El árbol de decisión se sitúa
en segundo lugar según AUC, exactitud y F1, y el método de vecinos en tercero.

Tres observaciones merecen atención.

La primera concierne a la referencia trivial. Su exactitud de 0.741 queda apenas seis centésimas
por debajo de la del método de vecinos, pese a que este último procesa cuarenta variables. Su AUC,
en cambio, es 0.500 frente a 0.835: la distancia entre ambos modelos resulta evidente en una
métrica y casi imperceptible en la otra.

La segunda se refiere al recall. Ningún modelo supera 0.52. El mejor de ellos identifica
aproximadamente la mitad de los pacientes que fallecen, cuestión que se analiza en la
Sección \ref{sec:recall}.

La tercera corresponde al orden entre el árbol y el método de vecinos, que se invierte según la
métrica: la precisión de vecinos, 0.721, supera a la del árbol, 0.701, mientras su recall es
notablemente inferior. El método de vecinos emite menos señalamientos y acierta en una proporción
mayor de ellos, a costa de omitir más casos.

\subsection{Curvas ROC}

La Figura \ref{fig:roc} presenta las curvas ROC de los tres clasificadores. Cada punto de la curva
corresponde a un umbral de decisión distinto, de modo que la representación describe el
comportamiento del modelo en todo el rango de umbrales y no únicamente en el valor predeterminado.

\begin{figure}[H]
\centering
\includegraphics[width=0.72\textwidth]{roc.png}
\caption{Curvas ROC sobre el conjunto de prueba. La diagonal corresponde a un clasificador
aleatorio.}
\label{fig:roc}
\end{figure}

La curva de la regresión logística domina a las restantes en prácticamente todo el recorrido. Las
de árbol y vecinos se entrecruzan, lo que anticipa el resultado de las pruebas de significancia
presentadas más adelante.

%============================================================================
\section{Matrices de confusión}
%============================================================================

Las métricas agregadas no indican qué tipo de error comete cada modelo. La Tabla
\ref{tab:confusion} descompone las clasificaciones en sus cuatro categorías.

\begin{table}[H]
\centering
\caption{Matrices de confusión sobre el conjunto de prueba. VN y VP son los aciertos; FP las falsas alarmas y FN los fallecimientos no detectados.}
\label{tab:confusion}
\small
\begin{tabular}{p{3.6cm} r r r r}
\toprule
\textbf{Modelo} & \textbf{VN} & \textbf{FP} & \textbf{FN} & \textbf{VP} \\
\midrule
Regresión logística & 1261 & 88 & 227 & 245 \\
Árbol de decisión & 1259 & 90 & 261 & 211 \\
Vecinos más cercanos & 1284 & 65 & 304 & 168 \\
\bottomrule
\end{tabular}
\end{table}


La columna de falsos negativos resulta la más relevante: corresponde a pacientes que fallecieron y
que el modelo no señaló. La regresión logística registra 227 casos, el árbol 261 y el método de
vecinos 304, sobre un total de 472 fallecimientos en el conjunto de prueba.

La comparación entre el árbol y el método de vecinos ilustra la naturaleza del compromiso. Este
último produce menos falsas alarmas, 65 frente a 90, pero omite 43 fallecimientos adicionales. Es
un modelo más conservador en sus señalamientos, propiedad que puede resultar deseable o no según
el costo relativo que se asigne a cada tipo de error.

%============================================================================
\section{Estabilidad de los resultados}
%============================================================================

Los valores presentados hasta aquí provienen de una única partición entre entrenamiento y prueba.
Una partición distinta produciría valores distintos, y conviene determinar la magnitud de esa
variación antes de atribuir significado a las diferencias observadas.

Para estimarla se aplicó validación cruzada estratificada de diez particiones sobre el conjunto de
entrenamiento. El procedimiento entrena y evalúa diez veces, de modo que proporciona una
distribución de valores en lugar de un único resultado.

\begin{table}[H]
\centering
\caption{AUC obtenido por cada modelo en las diez particiones de validación cruzada}
\label{tab:particiones}
\small
\begin{tabular}{l r r r}
\toprule
\textbf{Partición} & \textbf{Regresión logística} & \textbf{Árbol de decisión} & \textbf{Vecinos más cercanos} \\
\midrule
1 & 0.868 & 0.826 & 0.830 \\
2 & 0.874 & 0.826 & 0.824 \\
3 & 0.881 & 0.843 & 0.828 \\
4 & 0.867 & 0.841 & 0.847 \\
5 & 0.837 & 0.807 & 0.816 \\
6 & 0.844 & 0.813 & 0.810 \\
7 & 0.841 & 0.815 & 0.809 \\
8 & 0.842 & 0.802 & 0.809 \\
9 & 0.837 & 0.795 & 0.802 \\
10 & 0.847 & 0.807 & 0.798 \\
\midrule
Promedio & \textbf{0.854} & \textbf{0.817} & \textbf{0.817} \\
Desviación & 0.017 & 0.016 & 0.015 \\
\bottomrule
\end{tabular}
\end{table}


La regresión logística obtiene el mayor AUC en las diez particiones, sin excepción. Las
desviaciones se sitúan entre 0.015 y 0.017 en los tres modelos, magnitud superior a la diferencia
de 0.012 que separa al árbol del método de vecinos sobre el conjunto de prueba. En las propias
particiones, los promedios de ambos difieren en 0.000043, de modo que no resultan distinguibles a
partir de estos valores.

Los promedios de la validación cruzada resultan inferiores a los obtenidos sobre el conjunto de
prueba: 0.854 frente a 0.879 en la regresión logística. La diferencia corresponde a la variación
entre particiones, que es precisamente lo que este procedimiento permite estimar. Conviene por
tanto leer el valor de 0.879 como el obtenido sobre una partición concreta y no como la
expectativa del modelo.

%============================================================================
\section{Pruebas de significancia estadística}
%============================================================================

\subsection{Fundamento}

Las diferencias observadas entre modelos podrían originarse en el azar de la partición. Determinar
si corresponden a diferencias reales requiere pruebas estadísticas.

Se emplearon pruebas no paramétricas, que no presuponen que las métricas sigan una distribución
normal. En lugar de comparar promedios, comparan rangos —qué modelo quedó primero, segundo y
tercero en cada medición— y signos —cuántas veces uno superó a otro—. Con diez mediciones, una prueba de
normalidad tendría poca potencia, de modo que se prefirió no asumirla y emplear pruebas basadas en
rangos y signos.

El procedimiento consta de dos partes, aplicadas sobre conjuntos distintos.

\subsection{Friedman y Wilcoxon sobre validación cruzada}

Los tres modelos se evaluaron sobre las mismas diez particiones, lo que convierte las mediciones
en pareadas: cada partición constituye una prueba que los tres modelos rindieron en idénticas
condiciones.

La prueba de Friedman contrasta la hipótesis de que los tres modelos presentan el mismo
desempeño. Su resultado fue $\chi^2 = 15.000$ con $p = 0.0006$\unskip, de modo que la hipótesis se
rechaza: existen diferencias entre los modelos.

Establecida la existencia de diferencias, la prueba de Wilcoxon de rangos con signo determina
entre qué pares se verifican. El orden importa: aplicar comparaciones por pares sin contrastar
previamente la hipótesis global incrementa la probabilidad de detectar diferencias inexistentes.

Dado que se realizan tres comparaciones, los valores $p$ se multiplican por tres conforme a la
corrección de Bonferroni. Cuantas más comparaciones se efectúan, mayor es la probabilidad de que
alguna resulte significativa por azar, y la corrección compensa ese efecto.

\subsection{McNemar sobre el conjunto de prueba}

Las diez particiones comparten registros de entrenamiento entre sí, por lo que no son
estrictamente independientes y sus valores $p$ resultan optimistas. Esta limitación motiva una
segunda prueba sobre datos que ningún modelo utilizó.

La prueba de McNemar compara dos clasificadores paciente por paciente sobre el conjunto de prueba.
Considera únicamente los casos en que ambos modelos difieren: aquellos en que uno acierta y el
otro falla. Bajo la hipótesis de equivalencia, esos casos deberían repartirse por igual entre
ambos modelos. Se empleó la versión exacta, basada en la distribución binomial.

\subsection{Resultados}

La Tabla \ref{tab:pruebas} presenta ambas pruebas para los tres pares de modelos.

\begin{table}[H]
\centering
\caption{Comparación por pares. Los valores $p$ incluyen corrección de Bonferroni para tres comparaciones; el asterisco indica significancia al 5\%.}
\label{tab:pruebas}
\small
\begin{tabular}{p{5.2cm} r c r}
\toprule
\textbf{Comparación} & \textbf{Wilcoxon} & \textbf{Aciertos exclusivos} & \textbf{McNemar} \\
\midrule
Regresión logística vs. Árbol de decisión & 0.0059* & 95 / 59 & 0.0139* \\
Regresión logística vs. Vecinos más cercanos & 0.0059* & 123 / 69 & 0.0004* \\
Árbol de decisión vs. Vecinos más cercanos & 1.0000 & 118 / 100 & 0.7485 \\
\bottomrule
\end{tabular}
\end{table}


Los dos procedimientos coinciden. La regresión logística difiere significativamente del árbol de
decisión y del método de vecinos, tanto en validación cruzada como sobre el conjunto de prueba.
El árbol y el método de vecinos, en cambio, no presentan diferencias significativas por ninguno de
los dos procedimientos: $p = 1.0000$ en Wilcoxon y $p = 0.7485$ en McNemar.

Este último resultado constituye el principal hallazgo de la evaluación. El árbol obtiene un AUC
de 0.847 frente a 0.835 del método de vecinos, diferencia que la tabla de resultados presenta como
una ventaja. Las pruebas indican que dicha ventaja no se distingue del azar: sobre el conjunto de
prueba, 118 pacientes fueron clasificados correctamente sólo por el árbol y 100 sólo por el método
de vecinos, reparto compatible con la equivalencia entre ambos.

Corresponde señalar una particularidad de los valores obtenidos. Las dos comparaciones que
involucran a la regresión logística arrojan exactamente $p = 0.0059$ en la prueba de Wilcoxon. No
se trata de una coincidencia: con diez mediciones, el menor valor que la prueba puede producir es
$2/2^{10} = 0.00195$, que multiplicado por tres conforme a la corrección de Bonferroni resulta en
0.0059. Ambas comparaciones alcanzaron el mínimo posible, lo que equivale a afirmar que la
regresión logística superó a los otros dos modelos en la totalidad de las particiones.

%============================================================================
\section{El problema del recall}
\label{sec:recall}
%============================================================================

El mejor modelo identifica el 51.9\% de los pacientes que fallecen. Dado que detectar esos casos
constituye el propósito del clasificador, el resultado requiere análisis.

El resultado es consistente con el desbalance de clases: al representar los fallecimientos una
cuarta parte de los registros, el umbral predeterminado de 0.5 deja fuera a buena parte de la
clase minoritaria.

Una corrección consiste en ponderar las clases de forma inversa a su frecuencia, de modo que los
errores sobre la clase minoritaria pesen más durante el entrenamiento y el punto de operación se
desplace. La Tabla \ref{tab:balanceo}
presenta el efecto sobre la regresión logística.

\begin{table}[H]
\centering
\caption{Efecto de la ponderación de clases sobre la regresión logística}
\label{tab:balanceo}
\small
\begin{tabular}{p{5cm} r r r r}
\toprule
\textbf{Configuración} & \textbf{AUC} & \textbf{Exactitud} & \textbf{Precisión} & \textbf{Recall} \\
\midrule
Sin ponderación & 0.879 & 0.827 & 0.736 & 0.519 \\
Con ponderación balanceada & 0.879 & 0.791 & 0.569 & 0.790 \\
\bottomrule
\end{tabular}
\end{table}


El recall asciende de 0.519 a 0.790, mientras la precisión desciende de 0.736 a 0.569. El modelo
pasa a detectar cuatro de cada cinco fallecimientos, pero entre los pacientes que señala la
proporción de aciertos se reduce.

El AUC se mantiene en 0.879, con una variación inferior a dos diezmilésimas. El resultado es
coherente: la ponderación no altera de manera apreciable el ordenamiento de los pacientes por
riesgo, que es lo que el AUC mide, sino el umbral a partir del cual ese ordenamiento se traduce en
una decisión.

La elección entre ambas configuraciones no se resuelve mediante criterios estadísticos. Requiere
determinar el costo relativo de una falsa alarma frente al de un fallecimiento no detectado,
cuestión que corresponde al contexto de aplicación y no al modelo.

%============================================================================
\section{Conclusiones}
%============================================================================

La evaluación de los tres clasificadores arroja tres resultados.

El primero es que la regresión logística supera a los otros dos modelos, y que dicha superioridad
es estadísticamente significativa. Obtiene el mayor AUC en las diez particiones de validación
cruzada, y las pruebas de Wilcoxon y McNemar confirman la diferencia por ambos procedimientos.

El segundo es que el árbol de decisión y el método de vecinos no se distinguen entre sí. La
diferencia de doce milésimas de AUC que exhibe la tabla de resultados no sobrevive a las pruebas
de significancia. Presentarlos como modelos de desempeño diferente carecería de respaldo.

El tercero concierne al recall. Ningún modelo en su configuración predeterminada supera el 52\% de
detección sobre la clase de interés. La ponderación de clases eleva ese valor al 79\% a costa de
la precisión, compromiso cuya resolución depende del costo asignado a cada tipo de error.

La evaluación permite asimismo precisar el valor de cada métrica en este problema. La exactitud
resulta insuficiente, dado que un modelo que no aprende alcanza 0.741. El AUC describe
adecuadamente la capacidad de ordenamiento pero no informa sobre el comportamiento en el umbral de
decisión. La matriz de confusión y el par precisión–recall son los que permiten determinar qué
tipo de error comete cada modelo, que es la información necesaria para decidir cuál emplear.

\printbibliography[title={Referencias}]

\end{document}
